\begin{figure}[htbp]
    \centering
    \includegraphics[width=\textwidth]{figures/cultural_transmission_analysis.png}
    \caption{\textbf{Cultural Transmission of Bipedalism: Fire Circle Teaching Explains Feral Children and Socialized Apes.}
    \textbf{(A)} Genetic vs. cultural transmission dynamics over 100 generations. Genetic transmission (blue) maintains stable bipedalism proficiency near 100\%, while cultural transmission without teaching (red) shows rapid decay to feral baseline ($<$10\%) within 20 generations. Gray zones indicate cultural maintenance threshold (dashed) and feral failure zone (solid). Feral children fall into failure zone without fire circle teaching .
    \textbf{(B)} Socialized apes (Koko, Kanzi) achieve 80\% human-level bipedalism proficiency through systematic teaching (green), while wild apes without teaching (red dashed) remain at 20\% baseline .
    \textbf{(C)} Fire circle teaching network structure showing skill diffusion from expert teachers (red nodes, center) to learners (green gradient indicates skill level). Network topology enables rapid skill transmission through high-connectivity hub structure .
    \textbf{(D)} Population-level skill diffusion through teaching network over 20 time steps, showing mean skill level (blue line) increasing from 45\% to 85\% with teaching (shaded region = ±1 SD).
    \textbf{(E)} Population-level maintenance of bipedalism over 200 generations. With fire circle teaching (red), bipedalism remains stable at 5-10\% population frequency despite lack of genetic fixation. Without teaching (gray), trait is lost within 25 generations .
    \textbf{(F)} Teaching fidelity sensitivity analysis. Bipedalism skill level remains near zero until teaching fidelity exceeds 80\% threshold (dashed line), then jumps to stable maintenance at fire circle fidelity level (85-90\%, dotted line). This demonstrates critical threshold effect for cultural transmission .}
    \label{fig:cultural_transmission}
    \end{figure}


    \begin{figure}[htbp]
        \centering
        \includegraphics[width=\textwidth]{figures/fire_encounter_analysis.png}
        \caption{\textbf{Fire Encounter Inevitability in C4 Grassland Environments: Validation of Fire Circle Hypothesis.}
        \textbf{(A)} Baseline fire encounter probability calculation for typical hominid territory (10 km², 70\% C4 coverage, 150-day dry season). Using Equation 1 with conservative parameters yields P(Fire) = 1.0000, with expected 18.4 fires per season, establishing fire encounters as statistically inevitable .
        \textbf{(B)} Sensitivity analysis showing fire encounter probability remains $>$99.5\% (dark red) across realistic parameter ranges. Baseline parameters (white star) and variation analysis (second star) both yield near-certain fire exposure. Even with minimal territory (5 km²) and low C4 coverage (40\%), P(Fire) $>$ 99\% .
        \textbf{(C)} Temporal evolution from 8-3 Ma showing C4 grass coverage expansion (green dashed line) from 20\% to 70\%, with fire encounter probability (blue) reaching 99.6\% by 8 Ma and approaching 100\% by 3 Ma. Pink shading indicates period of inevitable fire exposure coinciding with early hominid evolution .
        \textbf{(D)} Spatial variation across 1,000 simulated territories showing fire encounter probability tightly clustered at 1.000 (mean = 1.000, 99\% threshold shown as dashed line). Zero territories fall below 99\% threshold, confirming statistical inevitability .
        \textbf{(E)} Expected fire encounters per season across territories (n=1,000) follows normal distribution with mean = 19.1 fires (dashed line). All territories experience 10-40 fires annually, establishing fire as dominant environmental feature $CITE_{10}$.
        \textbf{(F)} Poisson distribution of fire encounters with λ = 19.1 shows P(0 fires) = 0.000000, confirming zero probability of fire-free existence. Modal value = 19 fires per season $CITE_{11}$.}
        \label{fig:fire_encounter}
        \end{figure}


\begin{figure}[htbp]
\centering
\includegraphics[width=\textwidth]{figures/mechanotransduction_analysis.png}
\caption{\textbf{Mechanotransduction: Behavior Induces Morphology Without Bipedalism Genes.}
\textbf{(A)} Acute bone remodeling following Wolff's Law demonstrates rapid morphological response to mechanical loading. Bipedal behavior (red) induces 2× bone density increase within 6 months, while quadrupedal behavior (green) shows slower, linear response. No genetic changes required—pure mechanotransduction $CITE_{12}$.
\textbf{(B)} Multi-generational phenotypic expression shows bipedalism maintained at 25\% population frequency with teaching (red) vs. 8\% without teaching (orange) over 20 generations. Without any teaching (green), bipedalism disappears by generation 2. Demonstrates epigenetic inheritance through cultural transmission, not genetic fixation $CITE_{13}$.
\textbf{(C)} Maladaptation persistence over 300,000 generations. Cultural behavior (red) maintains bipedalism at 80-95\% despite genetic fitness (blue) remaining at 30-50\%, creating persistent maladaptation gap (purple dashed, 0.3-0.5). Cultural transmission is FASTER than genetic adaptation, preventing optimization. After $>$300k generations, maladaptation persists, contradicting traditional adaptive hypotheses $CITE_{14}$.
\textbf{(D)} Developmental trajectories showing teaching is essential. Socialized human children (blue) reach 95\% bipedalism proficiency by age 3 through fire circle teaching. Feral human children (gray dashed) plateau at 20\% (failure range). Socialized apes (green) achieve 60\% through teaching, exceeding feral humans. Human bipedalism range (blue shading) is only accessible via cultural transmission $CITE_{15}$.
\textbf{(E)} Mechanical loading patterns for spine, knee, and hip. Bipedal posture (red) shows 1.0× spine loading, 1.2× knee loading, 0.9× hip loading. Quadrupedal (green) shows 0.6× spine, 0.5× knee, 0.5× hip. Bipedal knee experiences 2.4× higher loading than quadrupedal, explaining persistent knee pathology as maladaptation $CITE_{16}$.}
\label{fig:mechanotransduction}
\end{figure}








    \end{figure}
    \begin{figure}[htbp]
        \centering
        \includegraphics[width=\textwidth]{figures/posture_transition_analysis.png}
        \caption{\textbf{Horizontal→Vertical Transition: 90° Rotation of Extended Spine.}
        \textbf{(A)} Spine configurations showing muscle activation patterns. Horizontal sleeping (blue, T=0.10) shows minimal activation across all muscle groups. Vertical standing (red, T=0.30) shows moderate activation after rotation. Quadrupedal (black, T=0.80) shows high constant tension. Fire circle hypothesis: horizontal→vertical is simple rotation; traditional theory: quadrupedal→vertical requires complete restructuring $CITE_{24}$.
        \textbf{(B)} Horizontal sleeping provides 8-12 hours of extended spine calibration with ZERO TENSION (shaded region). All muscle groups (erector spinae, abdominals, hip flexors, gluteals) show minimal activation (0.02-0.10), establishing neurological baseline for extended posture. This creates the foundation for vertical standing $CITE_{25}$.
        \textbf{(C)} Vertical standing after rotation shows moderate muscle activation (0.26-0.36) with characteristic U-shaped pattern during standing bout (2 hours). Activation drops to minimum at 1 hour (extended posture zone), then increases. Pattern reflects rotated horizontal calibration $CITE_{26}$.
        \textbf{(D)} Quadrupedal posture shows high constant tension (0.60-0.63) across all muscle groups with continuous linear increase over 2 hours. No extended posture zone exists. Transition from this state to vertical would require complete neuromuscular restructuring $CITE_{27}$.
        \textbf{(E)} Energy cost comparison. Horizontal sleeping = 1.60 units (baseline). Vertical standing = 2.49 units (1.56× increase). Quadrupedal = 4.89 units (3.06× increase). Horizontal→vertical transition costs 0.89 units; quadrupedal→vertical would cost 2.40 units (2.7× more expensive) $CITE_{28}$.
        \textbf{(F)} Transition cost comparison. Fire circle pathway (horizontal→vertical, green) requires only rotation (cost = 0.20). Traditional pathway (quadrupedal→vertical, red) requires complete restructuring (cost = 0.90, 4.5× higher). Fire circle hypothesis explains lower transition barrier $CITE_{29}$.
        \textbf{(G)} Learning curves. Vertical standing (red) requires extensive learning (cultural transmission), reaching 100\% stability after 80 practice trials. Quadrupedal (green dashed) is innate (95\% stability from birth). Yellow box emphasizes vertical requires LEARNING, supporting cultural transmission hypothesis $CITE_{30}$.
        \textbf{(H)} Standing duration capacity increases from 2 minutes to 60 minutes over 100 practice trials for vertical posture (red), while quadrupedal (green dashed) maintains constant 60-minute capacity. Demonstrates skill acquisition through teaching $CITE_{31}$.}
        \label{fig:posture_transition}
        \end{figure}
